% === §5.1 Marco conceptual y teórico ===
% Redactado por el Investigador — Fase 4

El constructo de \emph{aprendizaje profundo} constituye el eje conceptual vertebrador
de esta propuesta. En la tradición fenomenográfica inaugurada por Marton y Säljö
(1976) y posteriormente operacionalizada por Biggs (1999) en su modelo SOLO
(\emph{Structure of the Observed Learning Outcome}), el aprendizaje profundo se
distingue del procesamiento superficial por la intención del aprendiz de comprender
—no meramente de reproducir—, por la búsqueda activa de relaciones entre conceptos,
por la aplicación del conocimiento a problemas no estructurados y por la capacidad de
argumentar técnicamente sobre las soluciones propuestas \cite{marton1976qualitative,
biggs1999teaching}. En el dominio de la ingeniería, la presente propuesta
operacionaliza este constructo en tres dimensiones observables, alineadas con la
literatura contemporánea sobre educación en STEM \cite{freeman2014active,
theobald2020active}: (i) \emph{motivación intrínseca}, entendida como la disposición
autónoma del estudiante a involucrarse en tareas ingenieriles desafiantes más allá de
la recompensa extrínseca; (ii) \emph{argumentación técnica}, definida como la
capacidad de construir, justificar y comunicar razonamientos ingenieriles utilizando
evidencia, modelos y criterios disciplinares; y (iii) \emph{resolución de problemas
auténticos}, concebida como la transferencia de conocimientos teóricos a situaciones
abiertas, contextualizadas y semejantes a las que enfrenta el ingeniero en el sector
productivo.

Sobre este fundamento pedagógico se erige el andamiaje tecnológico de la propuesta,
articulado en torno al paradigma emergente de la \emph{inteligencia artificial
agéntica}. A diferencia de los sistemas de IA convencionales —reactivos, de
turno único y sin estado—, un agente autónomo es una entidad computacional que
percibe su entorno, razona sobre objetivos, planifica secuencias de acciones, ejecuta
herramientas externas y aprende de la retroalimentación \cite{russell2021artificial}.
Cuando múltiples agentes especializados colaboran bajo un protocolo de orquestación,
se configura una \emph{arquitectura multiagente} cuya capacidad colectiva supera la
de cualquier agente individual \cite{wang2024survey, xi2023rise}. En el contexto de
esta propuesta, el \emph{ecosistema de agentes autónomos} se define como una
arquitectura multiagente pedagógicamente fundada, compuesta por agentes con roles
didácticos diferenciados —tutor, evaluador, generador de problemas, simulador y
asistente de programación— que comparten un modelo del aprendiz, una memoria
pedagógica común y un mecanismo de \emph{scaffolding} adaptativo que gradúa la
asistencia en función de la competencia diagnosticada. La orquestación entre agentes
y herramientas externas se realiza mediante el \emph{Model Context Protocol} (MCP),
un estándar abierto para la interacción segura, interoperable y auditable entre
modelos de lenguaje, agentes y fuentes de datos contextuales
\cite{anthropic2024mcp}.

El \emph{modelado del aprendiz} —\emph{learner modeling}— constituye el puente
teórico y computacional entre las dos tradiciones anteriores. Originado en la
investigación sobre sistemas tutores inteligentes (ITS), el modelo del aprendiz es
una representación computacional del estado de conocimiento, las habilidades, las
estrategias cognitivas y los estados afectivos del estudiante, que el sistema utiliza
para adaptar la instrucción \cite{hooshyar2024systematic}. Frente a los modelos
tradicionales —basados en reglas de producción, redes bayesianas o \emph{knowledge
tracing}—, esta propuesta adopta un enfoque híbrido que combina procesamiento de
lenguaje natural (NLP) para el análisis semántico de producciones textuales y de
código, analítica del aprendizaje multimodal para la extracción de patrones de
interacción, y aprendizaje automático supervisado para la clasificación del nivel de
aprendizaje profundo en las tres dimensiones definidas. Este modelo opera como la
capa fundacional del ecosistema: alimenta al agente tutor con el perfil del aprendiz,
informa al módulo de \emph{scaffolding} adaptativo sobre el nivel de asistencia
requerido, y proporciona las variables independientes para la validación
cuasiexperimental del impacto.

El \emph{scaffolding adaptativo} —término acuñado por Wood, Bruner y Ross
(1976) a partir de la teoría sociocultural de Vygotsky (1978)— describe el soporte
temporal y contingente que un experto (humano o artificial) proporciona a un aprendiz
para que resuelva tareas que se sitúan en su \emph{zona de desarrollo próximo} (ZDP),
esto es, más allá de lo que puede hacer solo pero dentro de lo que puede lograr con
ayuda \cite{vygotsky1978mind, wood1976role}. La propuesta traduce este principio al
dominio computacional mediante un módulo que, a partir del nivel diagnosticado por el
modelo del aprendiz, regula tres parámetros de la interacción agéntica: la
profundidad de las explicaciones generadas, el grado de direccionalidad en la
retroalimentación y la complejidad de los problemas auténticos presentados. Este
mecanismo se complementa con la noción de \emph{comunidades de práctica} de Lave y
Wenger (1991), que fundamenta el diseño de las interacciones colaborativas entre
estudiantes mediadas por el ecosistema agéntico \cite{lave1991situated}, y con el
construccionismo de Papert y Harel (1991), que legitima la centralidad de la
construcción activa de artefactos —código, modelos de simulación, diseños de
circuitos— como vehículo privilegiado para el aprendizaje profundo en ingeniería
\cite{papert1991constructionism}.

Desde la perspectiva de la inteligencia artificial, la viabilidad del ecosistema
propuesto descansa en la convergencia de tres líneas tecnológicas maduras. En primer
lugar, los grandes modelos de lenguaje (LLMs) —tanto propietarios como abiertos—
proporcionan capacidades de generación de texto, código y razonamiento que, acopladas
con mecanismos de generación aumentada por recuperación (RAG) y 
técnicas de ingeniería de \emph{prompts}, permiten la producción de contenido
pedagógico contextualizado, la evaluación semántica de respuestas abiertas y la
 argumentación técnica en dominios ingenieriles \cite{park2023generative,
singhal2024large}. En segundo lugar, la madurez alcanzada por los protocolos de
orquestación agéntica (MCP, LangGraph, CrewAI) y por los entornos de ejecución de
código controlado (\emph{sandboxed execution}) hace factible la construcción de un
ecosistema multiagente robusto, seguro y auditables \cite{xi2023rise}. En tercer
lugar, la infraestructura de cómputo híbrida —que combina estaciones Apple Silicon
para inferencia local con servidores virtuales en nube (VPS) para orquestación y
analítica— garantiza la factibilidad operativa del despliegue en aulas universitarias
colombianas sin dependencia exclusiva de APIs comerciales externas.

Las bases teóricas descritas abordan de manera directa y diferenciada los tres
subproblemas planteados en la sección~\ref{sec:problematica}. Frente al SP1
—diagnóstico y modelado multidimensional del aprendiz—, la convergencia entre la
teoría del aprendizaje profundo (Marton y Säljö, Biggs), la analítica del aprendizaje
multimodal y el NLP proporciona el marco conceptual y las herramientas computacionales
para construir un modelo del aprendiz que capture las dimensiones de motivación,
argumentación y resolución de problemas auténticos con métricas trazables. Frente al
SP2 —arquitectura agéntica autónoma—, la teoría de agentes (Russell y Norvig), el
paradigma multiagente (Wang et al., Xi et al.), el MCP y los principios de
\emph{scaffolding} adaptativo (Vygotsky; Wood, Bruner y Ross) ofrecen el andamiaje
para diseñar un ecosistema que personalice trayectorias de aprendizaje a partir del
perfil del aprendiz. Frente al SP3 —validación empírica y transferencia tecnológica
TRL~6/7—, los marcos de evaluación de impacto educativo (Freeman et al., Theobald et
al.) y el propio enfoque TRL de ingeniería proporcionan los criterios de rigor para
demostrar causalidad y madurez tecnológica en la intervención. La integración de
estas bases en una arquitectura coherente constituye la novedad fundamental de la
propuesta.
